\section{Overall Agent Architecture}
Agent v2.0 is built on the \textbf{Orchestrator + Sub-agents + ReAct Loop} model. Instead of a fixed pipeline, Gemini 2.5 Flash acts as the Orchestrator: analyzing requests, planning, invoking appropriate tools, and synthesizing results.

The 5-step processing flow:
\begin{enumerate}
    \item \textbf{Intent Classification} --- Classify user intent
    \item \textbf{Clarification Gate} --- Ask for clarification if ambiguous
    \item \textbf{Memory Agent} --- Load user preferences from PostgreSQL
    \item \textbf{ReAct Loop} --- Plan $\rightarrow$ Execute $\rightarrow$ Observe cycle (max 8 iterations)
    \item \textbf{LLM Synthesis} --- Generate a natural language response
\end{enumerate}

\begin{figure}[H]
\centering
\resizebox{\linewidth}{!}{%
\begin{tikzpicture}[
    node distance=1.5cm and 1cm,
    process/.style={rectangle, draw=darkgreen, fill=darkgreen!10, rounded corners, text width=3.5cm, text centered, minimum height=1.2cm},
    model/.style={rectangle, draw=red!80!black, fill=red!10, rounded corners, text width=3.5cm, text centered, minimum height=1.2cm},
    db/.style={cylinder, draw=orange, fill=orange!20, shape border rotate=90, aspect=0.25, text width=2.5cm, text centered, minimum height=1.5cm},
    input/.style={rectangle, draw=blue, fill=blue!10, rounded corners, text width=3cm, text centered, minimum height=1cm},
    arrow/.style={thick,->,>=stealth}
]

% Nodes
\node[input] (query) {User Query};
\node[process, below left=of query, xshift=0.5cm] (bm25) {BM25 Keyword Search\\(label + color)};
\node[model, below right=of query, xshift=-0.5cm] (vector) {FashionSigLIP\\Vector Search (semantic)};

\node[db, left=of bm25] (bm25db) {BM25 Index};
\node[db, right=of vector] (qdrant) {Qdrant Vector DB};

\node[process, below=2.5cm of query] (rrf) {RRF Fusion \&\\Soft Filter};

\node[model, below=of rrf] (rerank) {BGE Reranker\\(multilingual)};
\node[process, below=of rerank] (synth) {Gemini 2.5 Flash\\Synthesis};

\node[db, left=of synth] (pg) {PostgreSQL\\(products)};
\node[input, below=of synth] (output) {Results \& Advice};

% Arrows
\draw[arrow] (query) -- (bm25);
\draw[arrow] (query) -- (vector);
\draw[arrow, dashed] (bm25db) -- (bm25);
\draw[arrow, dashed] (qdrant) -- (vector);
\draw[arrow] (bm25) -- (rrf);
\draw[arrow] (vector) -- (rrf);
\draw[arrow] (rrf) -- (rerank);
\draw[arrow] (rerank) -- (synth);
\draw[arrow, dashed] (pg) -- (synth);
\draw[arrow] (synth) -- (output);

\end{tikzpicture}
}% end resizebox
\caption{End-to-End Workflow Diagram of Fashion Agent v2.0}
\label{fig:workflow_v2}
\end{figure}

\section{Step 1 --- Intent Classification (Router)}

\begin{tcolorbox}[title=Method: LLM Few-shot Intent Classification]
Gemini classifies queries into \textbf{7 intent groups} with a confidence score. If confidence $< 0.6$, the agent switches to the clarification step instead of searching immediately.
\end{tcolorbox}

\begin{table}[H]
\centering
\begin{tabular}{p{2.9cm}p{4.8cm}p{4.8cm}}
\toprule
\textbf{Intent} & \textbf{Condition} & \textbf{Action}\\
\midrule
text\_search    & Find specific products         & Search immediately\\
outfit\_request & Request outfit suggestions     & Search + outfit advice\\
follow\_up      & Reference previous result      & Use session context\\
product\_select & Choose numbered items from recent results & Save selected product IDs\\
view\_selections & Ask to review saved picks      & Show selected items list\\
out\_of\_scope  & Unrelated to fashion           & Politely refuse\\
unclear         & Information too vague          & Ask for clarification\\
\bottomrule
\end{tabular}
\caption{Intent classification table with conditions and corresponding actions}
\end{table}

\begin{lstlisting}[caption={Intent Classification --- agent/intent\_classifier.py}]
@dataclass
class ExtractedSlots:
    category: Optional[str] = None
    color: Optional[str] = None
    fabric: Optional[str] = None
    fit: Optional[str] = None
    construction: Optional[str] = None
    aesthetic: Optional[str] = None
    selected_numbers: list[int] = field(default_factory=list)

@dataclass
class ClassifiedIntent:
    intent: str  # text_search|outfit_request|follow_up|product_select|view_selections|out_of_scope|unclear
    confidence: float
    filters: dict
    refined_query: str
    extracted_slots: ExtractedSlots = field(default_factory=ExtractedSlots)
    input_tokens: int = 0
    output_tokens: int = 0

def classify_intent(query: str, history: list | None = None) -> ClassifiedIntent:
    model = get_model()
    history_text = format_history_text(history, limit=4)
    full_prompt = INTENT_PROMPT.format(history_text=history_text) + query
    response = model.generate_content(full_prompt)

    usage = getattr(response, "usage_metadata", None)
    in_tok = getattr(usage, "prompt_token_count", 0) or 0
    out_tok = getattr(usage, "candidates_token_count", 0) or 0

    data = parse_llm_json(response.text)
    return ClassifiedIntent(
        intent=data.get("intent", "text_search"),
        confidence=float(data.get("confidence", 0.5)),
        filters=data.get("filters", {}),
        refined_query=data.get("refined_query", query),
        extracted_slots=_parse_slots(data),  # slot_* + selected_numbers
        input_tokens=in_tok,
        output_tokens=out_tok,
    )
\end{lstlisting}

\textbf{Explanation:}
\begin{itemize}
    \item \textbf{History-aware}: Uses the last 4 messages to understand context for follow-up queries.
    \item \textbf{Joint extraction}: One call returns intent, coarse filters, six detailed slots, and selected product numbers.
    \item \textbf{Token telemetry}: Prompt and output token counts are attached for logging and cost analysis.
    \item \textbf{Fallback}: On JSON parse error, the classifier defaults to \texttt{text\_search} with confidence 0.5.
\end{itemize}

\section{Step 2 --- Clarification Gate}

\begin{tcolorbox}[title=Method: Proactive Clarification]
When confidence $< 0.6$, the Agent does not guess --- it asks the user instead. This mechanism prevents incorrect searches, wastes fewer resources, and avoids returning irrelevant results.
\end{tcolorbox}

\begin{table}[H]
\centering
\begin{tabular}{ll}
\toprule
\textbf{Ambiguous Query} & \textbf{Agent's Clarification Question}\\
\midrule
``I want a shirt''        & ``Do you want a dress shirt, T-shirt, or jacket?''\\
``something nice''        & ``What occasion are you shopping for?''\\
``party outfit''          & ``Indoor or outdoor? What color do you prefer?''\\
``this style''            & ``Could you describe it in more detail?''\\
\bottomrule
\end{tabular}
\caption{Examples of clarification questions generated by the Clarification Gate}
\end{table}

\begin{lstlisting}[caption={Clarification Gate --- agent/clarification\_gate.py}]
@dataclass
class ClarificationResult:
    needs_clarification: bool
    question: str = ""

def check_clarification(query: str, history: list | None = None) -> ClarificationResult:
    """Called only when intent is unclear / low confidence."""
    lang = detect_language(query)
    fallback_q = FALLBACK_QUESTION.get(lang, FALLBACK_QUESTION["en"])

    try:
        model = get_model()
    except RuntimeError:
        return ClarificationResult(needs_clarification=True, question=fallback_q)

    history_text = format_history_text(history, limit=4)
    prompt = CLARIFICATION_PROMPT.format(query=query, history_text=history_text)

    try:
        response = model.generate_content(prompt)
        data = parse_llm_json(response.text)
        question = data.get("question", "")
        if question:
            return ClarificationResult(
                needs_clarification=True,
                question=question,
            )
    except Exception as exc:
        logging.getLogger(__name__).warning("Clarification LLM failed: %s", exc)

    return ClarificationResult(needs_clarification=True, question=fallback_q)
\end{lstlisting}

\section{Step 3 --- Memory Agent (PostgreSQL)}

\begin{lstlisting}[caption={Memory Agent --- agent/memory.py}]
def log_query(session_id, query, intent, filters):
    """Record query history to PostgreSQL JSONB."""
    entry = {
        "query": query, "intent": intent,
        "filters": filters, "ts": datetime.utcnow().isoformat()
    }
    with _get_conn() as conn:
        with conn.cursor() as cur:
            cur.execute("""
                UPDATE sessions
                SET query_history = query_history || %s::jsonb
                WHERE session_id = %s
            """, (json.dumps([entry]), session_id))
        conn.commit()

def add_liked_item(session_id, image_id, label, color):
    """Save a product the user has liked."""
    item = {"image_id": image_id,
            "label": label, "color": color}
    with _get_conn() as conn:
        with conn.cursor() as cur:
            cur.execute("""
                UPDATE sessions
                SET liked_items = liked_items || %s::jsonb
                WHERE session_id = %s
            """, (json.dumps([item]), session_id))
        conn.commit()
\end{lstlisting}

\textbf{Explanation:}
\begin{itemize}
    \item \textbf{JSONB append}: Uses PostgreSQL JSONB concatenation (\texttt{||}) to add new entries.
    \item \textbf{Accumulated preferences}: Automatically aggregates preferred categories and colors from history.
    \item \textbf{Limits}: Retains a maximum of 20 query history entries and 50 liked items.
\end{itemize}

\section{Step 4 --- The ReAct Loop}

\begin{tcolorbox}[title=Method: ReAct (Reason + Act) Loop]
ReAct is the central loop: Thought $\rightarrow$ Action $\rightarrow$ Observation $\rightarrow$ (repeat if needed, up to 8 iterations). Gemini both reasons and proactively selects the appropriate tool based on each step's observation.
\end{tcolorbox}

\begin{lstlisting}[caption={ReAct Loop --- agent/fashion\_agent.py}]
MAX_REACT_ITERATIONS = 8
LOW_CONFIDENCE_THRESHOLD = 0.5

def chat(query, session_id=None, api_key=None):
    """Main agent entry point - ReAct orchestrator."""
    # ... (intent, clarify, memory setup) ...

    search_query = intent_result.refined_query or query
    all_products = []
    observations = []
    threshold = LOW_CONFIDENCE_THRESHOLD

    for iteration in range(1, MAX_REACT_ITERATIONS + 1):
        # THOUGHT: Planner decides which tool to call
        plan = _plan(
            query=search_query,
            intent=intent_result.intent,
            preferences=preferences,
            observations=observations,
            iteration=iteration,
            max_iter=MAX_REACT_ITERATIONS)

        if not plan:  # Planner: "enough results already"
            break

        # ACTION: Execute at most 2 tools per iteration
        for tool_call in plan[:2]:
            tool_name = tool_call.get("tool", "search")
            tool_args = tool_call.get("args", {})
            obs, products = _execute_tool(
                tool_name, tool_args, preferences)
            observations.append(obs)
            if products:
                all_products = products

        # OBSERVE: Check result quality
        if all_products and all_products[0].score > threshold:
            break  # Results are good enough

        # Lower threshold after 4 iterations
        if iteration == 4:
            threshold -= 0.2

    # SYNTHESIZE response
    answer, styling = _synthesize_response(
        query, all_products, history, intent, preferences)
    return AgentResponse(answer=answer, products=all_products,
                         styling_suggestion=styling)
\end{lstlisting}

\textbf{3 available tools in ReAct:}
\begin{enumerate}
    \item \texttt{search(query, top\_k)} --- Calls the full hybrid search pipeline
    \item \texttt{memory\_enrich(query, preferences)} --- Enriches the query with user preferences
    \item \texttt{outfit\_hints(occasion, style)} --- Gemini generates outfit pairing suggestions
\end{enumerate}

\textbf{Stopping conditions:}
\begin{itemize}
    \item Planner returns \texttt{[]} (sufficient results found)
    \item Top product score $> \text{threshold}$ (0.5, lowered to 0.3 after 4 iterations)
    \item Maximum of 8 iterations reached
\end{itemize}

\section{Step 5 --- Gender Filtering and LLM Synthesis}

Before finalizing the results, the system applies demographic-aware post-processing. A specific \texttt{\_filter\_by\_gender} function validates the retrieved products against the user's recorded gender in PostgreSQL, ensuring that (for example) female-specific clothing is not inappropriately recommended to male users.

\begin{lstlisting}[caption={Gender Filtering and LLM Synthesis --- agent/fashion\_agent.py}]
SYNTHESIS_PROMPT = """You are a helpful fashion assistant.
Based on search results and user query, provide a natural
response in the same language as the user's query.

User Gender Context: {gender_context}
User query: {query}
Search results: {products_text}
User preferences: {preferences_text}
Conversation history: {history_text}

Instructions:
1. Respond in the same language (Vietnamese or English)
2. Tailor styling advice considering the {gender_context}
3. Briefly describe top recommendations and why they match
4. Include styling suggestions for outfit_request intent
5. Keep response concise (2-4 sentences)

Respond with JSON:
{{"answer": "<response>",
  "styling_suggestion": "<optional tips>"}}"""

def _synthesize_response(query, products, history,
                         intent, preferences, gender_context):
    """Gemini synthesizes a natural language response with gender awareness."""
    
    # Filter products by gender before synthesis
    filtered_products = _filter_by_gender(products, gender_context)
    
    products_text = "\n".join(
        f"{i}. {p.label} | {p.color} | {p.caption[:80]}"
        for i, p in enumerate(filtered_products, 1))
        
    prompt = SYNTHESIS_PROMPT.format(
        gender_context=gender_context,
        query=query, products_text=products_text,
        preferences_text=format_prefs(preferences),
        history_text=format_history(history[-6:]))
    response = model.generate_content(prompt)
    data = json.loads(response.text)
    return data["answer"], data.get("styling_suggestion", "")
\end{lstlisting}

\textbf{Explanation:}
\begin{itemize}
    \item \textbf{Gender Context Injection}: By explicitly stating the \texttt{gender\_context} in the prompt, Gemini can tailor the phrasing and styling tips appropriately (e.g., suggesting a ``sharp tailored fit'' for men vs. a ``flattering silhouette'' for women).
    \item \textbf{A/B Testing Capabilities}: The \texttt{gender\_hint\_enabled} flag allows developers to easily toggle this feature for controlled experimentation and research.
\end{itemize}

% ===========================================================================
% CHAPTER 6: AGENT v2.0 --- PATH 2: IMAGE-TO-IMAGE SEARCH
% ===========================================================================
